\section{结论与未来工作}
\label{sec:concl}

\noindent
在模型即服务系统中，自动扩缩容是同时实现高有效吞吐和高硬件利用率的关键，
但现有自动扩缩容机制既慢又是 stop-the-world，
其带来的性能开销显著限制了实际效果。
本文首先表明，借助基于网络、面向模型的多播，
模型自动扩缩容的数据平面可以在少于 $O(1)$ 缓存的条件下做到快速；
随后我们进一步说明，借助面向模型的远程执行，
数据平面还可以做到在线。
结合这两项技术后，
与当前最先进的、分别支持和不支持自动扩缩容的服务系统相比，
{\sys} 最多可带来 94.1\,\% 的性能提升和 19.46\,\% 的资源利用率提升。
我们相信，这项工作展示了由自动扩缩容赋能的模型即服务系统所具备的潜力与可行性。

尽管 {\sys} 的快速在线自动扩缩容向现代弹性服务系统迈出了关键一步，
仍有若干挑战值得继续研究。
首先，在我们的研究中发现，
扩缩容策略——即何时扩缩、如何扩缩——同样会显著影响系统效率。
这类策略高度依赖工作负载特征，
我们将其留作未来工作。
其次，{\sys} 当前聚焦于实例级扩缩容，
而现代模型也可以通过改变实例内部的并行配置来扩缩，
例如对混合专家（MoE）模型中的专家进行扩缩。
原则上 {\sys} 也适用于这类场景，
但更细致的探索仍有待未来开展。
